\begin{thebibliography}{14}
	\bibitem{astrakhan} Китаева В. Д., Суматохина А. С. Программная генерация задач ЕГЭ по математике на языке JavaScript по темам «Текстовые задачи» и «Вектора» // 75-я Международная студенческая научно-техническая конференция, посвящённая 95-летию АИРХ-АТИРПИХ-АГТУ : материалы конференции (Астрахань, 14–19 апреля 2025 г.). – Астрахань : АГТУ, 2025. – С. 810–813.	
	\bibitem{chas-ege} Тренажёр «Час ЕГЭ» [Электронный ресурс]. – URL: https://math.vsu.ru/chas-ege/sh/katalog.html (дата обращения: 11.06.2026).
	\bibitem{fipi} Федеральный институт педагогических измерений : открытый банк заданий ЕГЭ [Электронный ресурс]. – URL: https://fipi.ru/ege/otkrytyy-bank-zadaniy-ege .
	\bibitem{posobie} Момот Е. А., Арахов Н. Д. Разработка и внедрение ПО для сбора статистики результатов подготовки к ЕГЭ по математике профильного уровня // Актуальные проблемы прикладной математики, информатики и механики. – 2021. – С. 1–2.
	\bibitem{egemath} Открытый банк задач ЕГЭ по математике. Профильный уровень [Электронный ресурс]. – URL: https://prof.mathege.ru/ .
	\bibitem{ege} Единый государственный экзамен [Электронный ресурс] // Википедия : свободная энциклопедия. – URL: https://ru.wikipedia.org/wiki/Единый\_государственный\_экзамен .
	\bibitem{sdamgia} Портал «Решу ЕГЭ» [Электронный ресурс]. – URL: https://ege.sdamgia.ru/problem?id=27074 (дата обращения: 11.06.2026).
	\bibitem{fowler} Фаулер М. Рефакторинг: улучшение существующего кода. – СПб. : Питер, 2003. – 432 с.
	\bibitem{martin} Мартин Р. Чистый код: создание, анализ и рефакторинг. – М. : ДМК Пресс, 2019. – 464 с.
	\bibitem{mcconnell} Макконнелл С. Совершенный код. – М. : Русская редакция, 2005. – 896 с.
	\bibitem{oppenheim} Оппенгейм Д., Вудс Д. Разработка веб-приложений с помощью JavaScript и Node.js. – СПб. : Питер, 2020. – 576 с.
	\bibitem{js-good-parts} Крокфорд Д. Как устроен и работает JavaScript. – СПб. : Питер, 2019. – 240 с.
	\bibitem{sbornik} Прототипы заданий № 1–14 ЕГЭ [Электронный ресурс] // Easy-Physic.ru. – 2015. – URL: https://www.easy-physic.ru/wp-content/uploads/2015/11/prototipy-zadanij-1-14-ege-.pdf .
	\bibitem{golubev2015} Голубев А. А., Спасская Т. А. Сборник заданий по математике : учебное пособие [Электронный ресурс]. – Тверь : Тверской государственный университет, 2015. – 160 с. – ISBN 978-5-7609-0976-3. – URL: https://www.elibrary.ru/item.asp?id=23321749 .
\end{thebibliography}